%% cover_letter.tex — C1: Physics-Constrained Adaptive Slope
%% IEEE Transactions on Automatic Control
\documentclass[11pt,a4paper]{letter}
\usepackage[T1]{fontenc}
\usepackage[utf8]{inputenc}
\usepackage[margin=2.5cm]{geometry}
\usepackage{microtype}

\signature{Anoop~V.~Eluvathingal\\
           Ph.D.\ Candidate, Dept.\ of Electrical Engineering\\
           Indian Institute of Technology Madras\\
           Chennai 600036, India\\[2pt]
           \texttt{arun.eluvathingal@iitm.ac.in}}

\address{Department of Electrical Engineering\\
         Indian Institute of Technology Madras\\
         Chennai 600036, India}

\begin{document}
\begin{letter}{Editor-in-Chief\\
               \textit{IEEE Transactions on Automatic Control}}

\opening{Dear Editor,}

We submit for your consideration the manuscript:

\begin{center}
  \textbf{Physics-Constrained Adaptive Slope for Differential
  Protection: A Projected-Gradient Law with Lyapunov Stability
  Guarantees}
\end{center}

\textbf{Control-theoretic content.}
This paper casts a power engineering problem --- adapting the slope
parameter $k$ in percentage-restrain differential protection --- as a
formally stated online parameter estimation problem and provides a
complete stability analysis via a quadratic Lyapunov function.
The central result, Theorem~1, proves uniform ultimate boundedness (UUB)
of the tracking error $\tilde{k} = k - k^*$ with an explicit,
computable ultimate bound $\varepsilon_\infty$.  Corollary~1 shows
that the physics-derived constraint on the feasible set (the SAMBPS
veto gate $\kappa_n < 30$) is \emph{necessary and sufficient} for
the constraint set to be non-empty and the UUB guarantee to hold.
These are clean, self-contained control-theoretic results that use
power engineering as the motivating application but stand independently.

\textbf{Novelty over prior adaptive relaying work.}
Prior adaptive differential relaying work (\textit{e.g.}, Kasztenny
2000, Sachdev 1989, Phadke 1988) relies on offline tuning or MRAS
frameworks proven only for LTI plants.  The proposed law handles a
\emph{time-varying} unknown $k^*(t)$ under a nonlinear
physics-derived constraint, and provides a closed-form ultimate bound
that is verifiable at design time.

\textbf{Five contributions.}
\begin{enumerate}
  \item Formal problem statement: slope as constrained time-varying
        parameter with a physics-derived feasible set $\mathcal{K}(\kappa_n)$.
  \item Projected-gradient adaptive law (Algorithm~1) implementable
        at 100\,ms cycle rates on embedded DSPs.
  \item Lyapunov UUB theorem with explicit bound
        $\varepsilon_\infty$ and convergence rate $\alpha$.
  \item Necessity-and-sufficiency proof linking the control-theoretic
        result to the model-based protection veto gate.
  \item $7\times$ RMSE improvement over fixed-slope alternatives on a
        33\,kV mixed-IBR feeder (2400 simulation scenarios).
\end{enumerate}

\textbf{Fit to the journal.}
\textit{IEEE Transactions on Automatic Control} has a strong tradition
of Lyapunov-based adaptive control theory and its application to
engineering systems.  This paper makes a new connection between
adaptive parameter estimation with projection and the power system
protection domain, where formal stability guarantees are currently
absent from the literature.

\textbf{Suggested reviewers.}
\begin{itemize}
  \item Prof.\ Petros Ioannou (University of Southern California)
        --- robust adaptive control, projection operators
  \item Prof.\ Romeo Ortega (Universit\'{e} Paris-Saclay)
        --- Lyapunov stability, port-Hamiltonian methods
  \item Prof.\ Anna Stefanopoulou (University of Michigan)
        --- constrained adaptive estimation
  \item Prof.\ Anuradha Annaswamy (MIT) --- adaptive control theory
\end{itemize}

This manuscript is not under review elsewhere.
All authors have approved the submission.

\closing{Sincerely,}
\end{letter}
\end{document}
